% 03_psychological_exposure.tex

\section{Primitive Foundation for Reversed Exposure Costs}
\label{sec:psychological-exposure}

The model in Section~\ref{sec:primal-model} needs one economic input:
\[
\Phi_H>\Phi_L.
\]
This is the reversed exposure-cost ranking. In the standard Spence ordering, high types can bear ambitious signals because their production cost is lower. Here the relevant cost is not production effort but falsification exposure: an ambitious signal creates a claim that can be adversely repriced if high scrutiny arrives. High types may then have more to lose from carrying exposed ambition into the future.

The KL construction below is only a primitive foundation for this reversed ranking. It shows that \(\Phi_H>\Phi_L\) can come from primitive self-accuracy levels, not from an equilibrium history, accumulated reputation capital, or the receiver's later off-path attribution. The main model can therefore treat \(\Phi_\theta\) as a true-type exposure cost and let the D1 argument operate separately on beliefs about deviations.

Two distinctions matter. First, adverse repricing is not truth verification. A high-precision audit would tend to confirm genuinely accurate high types; that is not the technology studied here. Second, the exposure term remains attached to the sender's true type. The receiver's posterior affects the current reputational return to \(s_A\), but it does not redefine \(\Phi_\theta\).

\subsection{Reduced-form exposure payoff}

Let $j\in\{L,H\}$ denote the current judgment-intensity state, with associated scrutiny multiplier $\mu_j>0$. The low state $L$ represents ordinary evaluation and the high state $H$ represents intensified evaluation. The state variable $j$ is not a political primitive. It is an exogenous state of the social evaluation technology.

After $s_A$, observers form beliefs not only about the sender's ability $\theta$, but also about the accuracy of the sender's self-model. Write
\[
m\in\{0,1\}
\]
for the event that the sender's self-model is accurate, where $m=1$ denotes accurate self-assessment. The sender cares not only about what observers infer about $\theta$, but also about what observers infer about the sender's own assessment of $\theta$.

For a finite hierarchy of beliefs $(\mu^2,\ldots,\mu^k)$, let
\[
L(\mu^2,\ldots,\mu^k;j,x)
\]
denote the epistemic penalty in judgment state $j$ and exposure state $x\in\{0,1\}$. The exposure state records whether the sender previously chose the ambitious signal. Choosing $s_A$ sets $x=1$ next period; choosing $s_P$ clears exposure next period.

The augmented one-period payoff can be written as
\[
u_\theta(s,j,x;\mu^2,\ldots,\mu^k)
=
v_\theta(s)-c_\theta(s)-\mu_j L(\mu^2,\ldots,\mu^k;j,x),
\]
where $v_\theta(s)$ is the ordinary reputational or material payoff from the signal, $c_\theta(s)$ is the intrinsic cost of the signal, and the final term is the judgment-weighted exposure penalty.

The next assumption is purely regularity. It rules out discontinuities inside a fixed judgment state. The discontinuities studied below will come from changes in the transition probability into high scrutiny, not from discontinuous belief-dependent utility within a regime.

\begin{assumption}[Regime-conditional continuity]
\label{ass:belief-continuity}
For each fixed $j\in\{L,H\}$ and $x\in\{0,1\}$, the function
\[
L(\mu^2,\ldots,\mu^k;j,x)
\]
is continuous in the relevant finite hierarchy of higher-order beliefs. The action set, type space, and exposure state space are finite.
\end{assumption}

\begin{lemma}[Regime-conditional continuity]
\label{lem:belief-utility-continuity}
Under Assumption~\ref{ass:belief-continuity}, for each fixed judgment-intensity state $j$, the induced belief-dependent payoff is continuous in the relevant higher-order belief hierarchy. On any open region where the relevant incentive constraints are strict, the associated pure-strategy equilibrium selection is locally constant.
\end{lemma}

\begin{proof}
Fix $j$ and $x$. The ordinary payoff terms $v_\theta(s)$ and $c_\theta(s)$ are defined on a finite action and type space. The multiplier $\mu_j$ is constant within the state $j$. By Assumption~\ref{ass:belief-continuity}, the exposure penalty $L(\mu^2,\ldots,\mu^k;j,x)$ is continuous in the relevant belief hierarchy. Multiplication by the fixed scalar $\mu_j$ and addition of finite ordinary payoff terms preserve continuity. Thus, the regime-conditional payoff is continuous.

Since the type and action spaces are finite, strict incentive inequalities are preserved under sufficiently small perturbations of beliefs and payoff parameters within a fixed state. Therefore, wherever the relevant incentive constraints are strict, the same pure-strategy profile remains selected in a neighborhood of the primitive. Any discontinuous change in equilibrium behavior must therefore come from a change in the relevant state, transition probability, or binding incentive constraint, rather than from a discontinuity of belief-dependent utility within the fixed state.
\end{proof}

\subsection{Primitive self-accuracy levels}

The exposure index is built from a level channel, not a precision channel. Let
\[
p_\theta:=\Pr(m=1\mid \theta)
\]
denote the primitive probability that a type-$\theta$ sender has an accurate self-model. The key ordering is
\[
1-\delta>p_H>p_L>\bar p>\delta>0.
\]
Higher types are genuinely more self-accurate in the primitive environment. The lower benchmark $\bar p$ is the assessment induced by adverse repricing under high scrutiny.

This assumption should not be read as signal precision. Precision is a fidelity concept: a more precise audit reveals the true state more accurately. If high types are genuinely more accurate, a fidelity-improving audit tends to confirm them, not punish them. The model instead studies adverse social repricing. In a high-scrutiny state, an exposed self-presentation can be reinterpreted downward with a positive probability that is not driven by the sender's true merit. The economic loss is increasing in the amount of self-model reputation-at-risk carried by the true type.

Let $\alpha\in(0,1]$ be the probability, conditional on high scrutiny and exposure, that the ambitious self-presentation is adversely repriced. This probability is taken to be independent of type in the baseline. It can be normalized into $\mu_H$, but keeping it visible makes the interpretation clear: high scrutiny is not a deterministic truth audit. It is a state in which exposed self-presentations face a type-independent risk of adverse reinterpretation.

The epistemic exposure index is
\[
\Phi_\theta
=
\alpha\,
D_{\mathrm{KL}}
\left(
\operatorname{Bern}(p_\theta)
\,\middle\|\,
\operatorname{Bern}(\bar p)
\right).
\]
The direction of the divergence reflects the economic object. The loss is not the ordinary probability of failure. It is the collapse from a carried self-accuracy level to a lower adverse benchmark. Judgment intensity prices this exposure:
\[
C_E(\theta,j)=\mu_j\Phi_\theta.
\]

Thus $\Phi_\theta$ is neither history-induced reputation capital nor signal precision. It is the true-type-indexed loss from adverse repricing of self-model exposure. The receiver's later off-path attribution may affect the reputational return to a deviation, but it does not redefine $\Phi_\theta$.

Because the self-model prior enters this loss mechanically, higher prior self-accuracy creates a larger exposure loss when it is repriced against the lower benchmark. That is the only role of the KL calculation: it delivers a primitive sufficient condition for \(\Phi_H>\Phi_L\).

\subsection{Exposure-cost ranking}

\begin{lemma}[Primitive foundation for reversed exposure costs]
\label{lem:bayesian-surprise-ordering}
Suppose
\[
1-\delta>p_H>p_L>\bar p>\delta>0.
\]
Define
\[
\Phi_\theta
=
\alpha\,
D_{\mathrm{KL}}
\left(
\operatorname{Bern}(p_\theta)
\,\middle\|\,
\operatorname{Bern}(\bar p)
\right)
\]
with $\alpha>0$, and let
\[
C_E(\theta,j)=\mu_j\Phi_\theta.
\]
Then
\[
\Phi_H>\Phi_L.
\]
Consequently, the exposure cost has increasing differences in type and judgment intensity. For every $\mu'>\mu$,
\[
C_E(\theta_H,\mu')-C_E(\theta_H,\mu)
>
C_E(\theta_L,\mu')-C_E(\theta_L,\mu).
\]
Equivalently, the marginal cost of higher judgment intensity is larger for the high-exposure type.
\end{lemma}

\begin{proof}
For Bernoulli beliefs,
\[
D_{\mathrm{KL}}
\left(
\operatorname{Bern}(p)
\,\middle\|\,
\operatorname{Bern}(\bar p)
\right)
=
p\log\frac{p}{\bar p}
+
(1-p)\log\frac{1-p}{1-\bar p}.
\]
Differentiating with respect to $p$ gives
\[
\frac{\partial}{\partial p}
D_{\mathrm{KL}}
\left(
\operatorname{Bern}(p)
\,\middle\|\,
\operatorname{Bern}(\bar p)
\right)
=
\log\frac{p(1-\bar p)}{\bar p(1-p)}.
\]
This derivative is positive whenever $p>\bar p$. Since $p_H>p_L>\bar p$, it follows that
\[
\Phi_H>\Phi_L.
\]
Therefore, for every $\mu'>\mu$,
\[
\begin{aligned}
&\big[C_E(\theta_H,\mu')-C_E(\theta_H,\mu)\big]
-\big[C_E(\theta_L,\mu')-C_E(\theta_L,\mu)\big]\\
&\qquad
=(\mu'-\mu)(\Phi_H-\Phi_L)>0.
\end{aligned}
\]
Thus higher judgment intensity raises the exposure cost of the high-self-accuracy type more than that of the low-self-accuracy type.
\end{proof}

Lemma~\ref{lem:bayesian-surprise-ordering} is not a separate equilibrium mechanism. It only supplies one primitive route to the exposure-cost ranking used by the main signaling argument. In a Spence environment, high types are more willing to bear the ambitious signal because they face a lower ordinary cost of producing it. Here, high types can face a higher exposure cost because the signal creates more future falsification risk for them. The reversal is therefore economic: the relevant cost ranking changes when the cost is exposure rather than production effort.

\begin{remark}[Repricing, not precision]
The sign of Lemma~\ref{lem:bayesian-surprise-ordering} depends on the level channel. Replacing $p_\theta$ with signal precision would change the model. A higher-precision audit is a fidelity technology and tends to confirm an accurate high type. It therefore does not generate top-down withdrawal. The primitive needed here is a type-indexed self-accuracy level combined with an adverse repricing state, not a variance or precision parameter.
\end{remark}

\subsection{From cost ranking to signaling reversal}

The preceding lemma converts primitive self-accuracy into exposure. The main model only uses the resulting ranking \(\Phi_H>\Phi_L\). A high type has more to gain from separation, but also more to lose from a future state that adversely reinterprets her ambitious signal.

Because $\Phi_\theta$ is a true-type exposure cost, falsification cost is not homogeneous across types: the ranking between types can reverse under repricing.

Let $B_\theta$ denote the current net separation rent from choosing $s_A$ rather than $s_P$ in the low-judgment state, before future high-scrutiny exposure is priced. Let
\[
q=\Pr(j_{t+1}=H\mid j_t=L)
\]
be the transition probability into high scrutiny. If choosing $s_A$ creates the one-period shadow, the low-state payoff gain from choosing $s_A$ can be written as
\[
D_\theta^L(q)=B_\theta-\beta q\mu_H\Phi_\theta.
\]
The associated threshold is
\[
q_\theta^*=\frac{B_\theta}{\beta\mu_H\Phi_\theta},
\]
whenever $B_\theta>0$.

Under the maintained static-separation case \(B_H>0\) and \(B_L<0\), only the high type has a positive current rent from \(s_A\). The operative threshold is therefore the high-type IC threshold
\[
q_{\mathrm{IC}}:=q_H^*=\frac{B_H}{\beta\mu_H\Phi_H}.
\]
For \(q>q_{\mathrm{IC}}\), the high type withdraws even while \(s_A\) is still read as a high-type signal. Whether pooling on \(s_P\) then survives off-path resurrection is the D1 question taken up in Section~\ref{sec:belief-loop}. The reversal relative to standard Spence is that the high type's constraint is the one that binds first: the same primitive that makes ambitious separation valuable also gives the high type larger exposure, \(\Phi_H>\Phi_L\).
